\section{서론}

평지 속도와 불연속 지형 적응성은 서로 다른 접촉 기구를 요구한다. 바퀴는 구름
구속을 유지하고 발을 반복 배치하지 않아도 되지만, 관절형 다리는 접점을 선택하고
그 접점에 대해 몸체를 움직일 수 있다. 바퀴--다리 로봇은 보통 다리 끝에 별도의
구동 바퀴를 유지한다 \cite{tedeschi2017cassino,bjelonic2019keep}. 반대로 회전
다리 로봇은 적은 액추에이터로 뛰어난 내재적 이동성을 얻지만
\cite{saranli2001rhex,quinn2002whegs,zhang2023qwhex}, 접촉 경로와 몸체 운동이
강하게 결합된다.

본 연구의 질문은 다음과 같다. \emph{완전 관절형 6족 로봇의 마지막 구동 링크를
어떻게 형상화해야 같은 관절과 접촉면으로 발 배치, 구름 주행, 계단 모서리 결합을
모두 수행할 수 있는가?} 제안 플랫폼은 다리당 3개씩 총 18개 액추에이터와 열린
C자형 TPU 호형 말단을 갖는다(그림~\ref{fig:teaser-ko}). 모드 전환을 위한 기계적
변형은 없다. 대신 보행에서는 호의 최저 패치를 고정하고, 주행에서는 패치를 호를
따라 이동시키며, 계단에서는 호의 경계 또는 트레드가 계단 코에 반력을 전달한다.

\begin{figure*}[t]
  \centering
  \begin{minipage}{0.32\textwidth}\centering
    \includegraphics[width=\linewidth]{flat_articulated-walk.jpg}\\[-1mm]
    \small (a) 관절형 보행
  \end{minipage}\hfill
  \begin{minipage}{0.32\textwidth}\centering
    \includegraphics[width=\linewidth]{flat_full-roll.jpg}\\[-1mm]
    \small (b) 전신 구름 주행
  \end{minipage}\hfill
  \begin{minipage}{0.32\textwidth}\centering
    \includegraphics[width=\linewidth]{stairs_200mm_full-scone.jpg}\\[-1mm]
    \small (c) 200 mm 계단 시뮬레이션
  \end{minipage}
  \caption{하나의 구동 C자형 말단을 세 가지 접촉 상태로 재사용한다. 모든 영상은
  결과 절의 벤치마크 조건에서 재현 가능하게 캡처한 MuJoCo 장면이며, 렌더링
  자체는 실물 성능 근거가 아니다.}
  \label{fig:teaser-ko}
\end{figure*}

본 논문의 기여는 세 가지다. 첫째, 다기능 호형 말단 형태와 그 구름/계단 접촉의
명시적 기하 모델이며, 여기에는 회전만으로는 넘을 수 없는 단높이의 닫힌 형태
경계가 포함된다. 둘째, \gaitname{}이다. 이는 삼각보 보행 궤적과 연속 호 회전을
같은 여섯 개 말단 관절에서 \emph{동시에} 수행하는 제어기로, 속도 모드 회전
속도를 보행 궤적의 미분에 중첩한다. 아울러 매끄러움과 주기 평균 이송 비율을
닫힌 형태로 특성화한 위상 게이트 변형을 함께 제시한다. 셋째, 동일 모델에서
관절 운동만, 말단 회전만, 결합 제어를 비교하는 절제 벤치마크이며, 해석적 접촉
속도 모델이 구름 조건의 실측 미끄럼을 $5\%$ 이내로 예측한다.

증거의 범위는 여기서 한 번에 밝힌다. 본 논문의 모든 수치는 결정론적 MuJoCo
모델에서 나왔고, 평지·계단·전환의 공칭 조건은 단일 시행이며 불규칙 지형 조건은
3개 seed를 사용한다. 이는 하나의 모델에서 세 제어 체계를 분리하고 기하 경계를
검정하기에는 충분하지만, 신뢰구간이나 실물 성능 주장에는 충분하지 않다. 실물
시연 자료는 존재하나 아래의 모든 수치 서술에서 의도적으로 제외했다.
